\documentclass[11pt,a4paper]{article}
\usepackage[margin=22mm]{geometry}
\usepackage{amsmath,amssymb}
\usepackage[hidelinks]{hyperref}
\setlength{\parindent}{0pt}
\setlength{\parskip}{5pt}
\raggedright
\begin{document}
{\Large What count accuracy does and does not establish for own motion}\\[4pt]
8 October 2026 --- source/static qualification; no joined law adopted

\textbf{Scope.} Source-supported CM and conditional numerical checks, followed
by reviewer-derived finite-birth and objective limits. These are not a traffic
simulation, reached metering comparison, full endogenous theorem or gradient
certificate. The source audit distinguishes the supplied 2005 journal from
the nearest cited 2004 thesis and unverified geometry manuscript.

\section*{1. The feedback quantity needs a definite operator}
Use free speed $u>0$, backward-wave magnitude $w>0$, jam density $\kappa>0$,
and continuous density $\rho$:
\begin{equation}
 F(\rho)=\min\{u\rho,w(\kappa-\rho)\},\quad
 \rho_c=\frac{w\kappa}{u+w},\quad
 V(\rho)=\begin{cases}u,&0\le\rho\le\rho_c,\\
 w(\kappa/\rho-1),&\rho_c<\rho\le\kappa.\end{cases}
\end{equation}
Units: $u,w,V$ in m/s; $\rho,\kappa$ in equivalents/m; $F$ in equivalents/s.
On a fixed interior downstream support of length $\ell>0$, compare
\begin{equation}
 v_{\rm density}=V\!\left(\frac1\ell\int_x^{x+\ell}\rho\,dy\right),\quad
 v_{\rm speed}=\frac1\ell\int_x^{x+\ell}V(\rho)\,dy,\quad
 v_{\rm ratio}=\frac{\int_x^{x+\ell}F(\rho)\,dy}
 {\int_x^{x+\ell}\rho\,dy}.
\end{equation}
The last expression assumes positive mass and an instantaneous spatial
diagnostic. It is not automatically the source's preceding-period Edie
lane-choice estimate or a prescribed CM estimator.

For equal support halves, $u=30$, $w=5$, $\kappa=3/20$,
$\rho_1=1/100$ and $\rho_2=1/10$:
\begin{equation}
 F_1=3/10,\quad F_2=1/4,\quad
 v_{\rm density}=95/11,\qquad v_{\rm speed}=65/4,\qquad
 v_{\rm ratio}=5.
\end{equation}
Thus all three natural readings differ on the same own field. The example
qualifies operator meaning; it is not hidden-information or control evidence.

The supplied 2006 Appendix A gives a constrained realized speed, desired-speed
recurrence and motion of the form
\begin{equation}
 v_j=\min\{v^{\rm dwn}_j,v^{\rm des}_j\},\qquad
 v^{\rm des}_{j+1}=v_j+a(v_j,\phi_j)\Delta t,\qquad
 \phi_{j+1}=\phi_j+v_j\Delta t.
\end{equation}
The own downstream estimate is the missing specification, not an extra
population. D04 supplies a density-average numerical precedent, not the
complete adapted mandatory-node/source/encounter law.

\newpage
\section*{2. A count bound alone does not close motion feedback}
In a \emph{source-free prescribed-path} diagnostic, $N_x=-\rho$. Compare
two compatible count fields on a common support with
\begin{equation}
 \|N_h-N\|_\infty\le\varepsilon_h,\qquad
 \bar\rho=\frac{N(t,x)-N(t,x+\ell)}{\ell}.
\end{equation}
Endpoint subtraction gives
\begin{equation}
 |\bar\rho_h-\bar\rho|\le\frac{2\varepsilon_h}{\ell}.
\end{equation}
The triangular-FD speed function is globally Lipschitz on $[0,\kappa]$,
with an admissible constant $L_V=w\kappa/\rho_c^2$: it is constant below
$\rho_c$ and its congested derivative has magnitude $w\kappa/\rho^2$.
Consequently
\begin{equation}
 |V(\bar\rho_h)-V(\bar\rho)|\le
 \frac{2L_V\varepsilon_h}{\ell}.
\end{equation}
J01's piecewise prescribed-path theorem bounds count error by
$\varepsilon_h\le3P\kappa\Delta x$ for its stated data class. For a
mesh-dependent support $\ell=m\Delta x$, the resulting estimate is only
\begin{equation}
 |V(\bar\rho_h)-V(\bar\rho)|\le\frac{6P\kappa L_V}{m}.
\end{equation}
This bound does not tend to zero with the mesh. It \emph{does not prove}
that actual speed errors fail to converge. A fixed positive physical support
does make this particular bound vanish, but changes the support limit and
cannot be silently interpreted as immediately downstream motion. Variable
end supports, moving path errors and endogenous sources need additional work.
Prescribed-path count convergence is not a joined feedback theorem.

\section*{3. Check a direct stationary boundary transcription}
Let an existing boundary-born restriction have prescribed speed $v=u$,
$\phi(t)=ut$, with empty/free-compatible receiving road, free downstream,
and attempted incoming flow $0<d\le Q=uw\kappa/(u+w)$. Ordinary compatible
traffic has $F(\rho)\le u\rho$, so a zero-relative-passing condition at that
moving path does not bind. Over the first interval the ordinary boundary can
admit $d\Delta t$.

The specific transcription that holds the restriction at $x=0$ with fixed
boundary capacity zero for the entire interval instead admits zero:
\begin{equation}
 \Delta C_{\rm moving}=d\Delta t,\qquad
 \Delta C_{\rm held}=0,\qquad
 \Delta C_{\rm moving}-\Delta C_{\rm held}=d\Delta t>0.
\end{equation}
It is a finite-step count mismatch consistent with a vanishing mesh error,
not a source-wide impossibility result. It exposes activation qualification
for the candidate's mandatory-boundary transfer. No full-TST difference follows
just from moving waiting mass outside versus inside before any exit occurs.

\newpage
\section*{4. Finite population bounds births, not clearance}
Use the admitted finite, forward-directed cell graph: through and discretionary
moves advance one section; the finite ramp joins the mainline once. Assign
a remaining-step rank $r_a$ to each cell or external waiting location, with
every accepted advance reducing rank by at least one and $0\le r_a\le K$.
No cyclic or backward population movement is admitted in this deduction.

For continuous masses $m_a$, define the nonnegative potential
\begin{equation}
 P(t)=\sum_a r_a m_a(t).
\end{equation}
Let $A_\infty<\infty$ be total offered mass after initialization and $H(T)$
the cumulative mass counted across all advancing moves through time $T$.
Conservation and decreasing ranks imply
\begin{equation}
 H(T)\le P(0)+K A(T)\le P(0)+K A_\infty.
\end{equation}
The constants $K$ and $S$ refer to the declared finite grid. This is not a
uniform record-count bound under arbitrary mesh refinement.
Moving waves do not move population backward. A flow split is debited once
per actual advance, in accordance with the source IT update; it is not counted
as both a same-cell lateral move and a second longitudinal move.

Let $S$ be the finite number of descriptor-generating streams, threshold
$\delta>0$, initial progress $r_s^0\in[0,\delta)$, and accepted generating
transfer $\eta_s(T)$. E1's cumulative floor gives
\begin{equation}
 B_s(T)=\left\lfloor\frac{r_s^0+\eta_s(T)}{\delta}\right\rfloor,
 \qquad \sum_s\eta_s(T)\le H(T).
\end{equation}
Hence, including all initial remainders and retaining all distinct births,
\begin{equation}
 B_{\rm new}(T)=\sum_s B_s(T)
 \le\frac{P(0)+K A_\infty}{\delta}+S.
\end{equation}
No threshold integerizes traffic mass. The bound is a finite-geometry
feasibility deduction, not a numerical cost or record-budget requirement.
It does not prove that traffic or restrictions eventually leave; a stationary
blocking update could still fail to clear despite finitely many births.

\newpage
\section*{5. Finite-time performance and derivatives are separate claims}
Under the same offered arrivals and initial population, conservative
full-system inventory includes outside waiting, ramp and mainline:
\begin{equation}
 M(t)=M_0+A(t)-E(t),\qquad
 J_T=\int_0^T M(t)\,dt.
\end{equation}
$E$ counts all exits. If a conditional comparison establishes a uniform exit
count error $|E_h-E|\le\varepsilon$ on $[0,T]$, then
\begin{equation}
 |J_{T,h}-J_T|\le T\varepsilon.
\end{equation}
This says nothing about the missing tail beyond $T$. Full $J$ requires
clearance or a qualified integrable tail. Normalizing by common total offered
and initial mass rescales the bound but does not remove that obligation.

Even uniform objective convergence does not guarantee its gradients. For
an arbitrary smooth $f$ and $h>0$, a generic calculus example is
\begin{equation}
 f_h(\theta)=f(\theta)+h\sin(\theta/h),\quad
 |f_h-f|\le h,\quad
 f_h'(\theta)=f'(\theta)+\cos(\theta/h).
\end{equation}
This is not a claimed traffic formula. It demonstrates the logical limit of
promoting count/objective accuracy into derivative accuracy. Threshold births
and active-set changes introduce additional obligations in the actual task.

\textbf{Disposition.} The physical source family and a concrete numerical
feedback precedent are supported. The displayed source-derived candidate is
not frozen: activation, feedback support/phase, joined accepted sources and
simultaneous paths require a declared completion. Finite births, finite-time
count/TST checks and a causal interface do not establish full clearance,
nontrivial metering leverage, record necessity or useful gradients.
\end{document}
